\section{Discussion}
\label{sec:discussion}

\subsection{Where the deformation enters}
Three mechanisms put deformation into a split-spectrum screen, and the experiments separate them.
\begin{enumerate}
  \item \emph{Cycle corrections.} Rounding an unwrapped field against its own low-pass, or against the filtered
        first-pass estimates, reads the deformation's curvature as unwrapping errors. Each false cycle of the main
        band is 3.21\,rad of screen, and the false cycles follow the deformation. The result is a fixed fraction of
        it, $-0.045$ to $-0.048$ on our synthetics and $-0.06$ without an exclusion, whatever the slip or the
        geometry. This is the sign-flipped copy the GUNW screens show. It needs no unwrapping error at all: the passes
        create the errors they correct. They also act at data edges, where their one-sided low-pass lags any steep
        gradient (Appendix~\ref{app:isce3}).
  \item \emph{Unwrapping a field that carries deformation across a rupture that splits the scene.} The side band
        carries all of the step and the double difference 4.4\% of it. Once that share exceeds about $\pi$, the
        cross-fault integer is a guess. The guess is wrong by one cycle of $\Delta$ (73\,rad, 1.4\,m of LOS) on D141,
        whose 3.8\,m step puts 1.38 cycles into $\Delta$, and on our synthetics from 4--6\,m of slip on. Where the
        unwrapper can go around the rupture's ends, this does not happen.
  \item \emph{Band co-location.} A side band sampled $\delta$ down range adds $12.1\,\delta$ times the deformation's
        range gradient. For strike-slip ruptures that a polar-orbiting SAR sees, the gradient mostly runs along
        azimuth, and ISCE3's 9\,m offset is worth well under a millimetre beyond 2\,km. Close to the trace, or for a
        rupture whose gradient points down range, it is centimetres.
\end{enumerate}
The reordered estimator removes the first two by construction, because its unwrapped field carries no deformation,
and the third by summing both bands over identical cells. The Antarctic pairs (Section~\ref{sec:antarctica}) show the
first two, and the products' main-band integers below, without an earthquake. Across a glacier's shear margins their
screen is ours plus:
\begin{itemize}
  \item one cycle of $\Delta$ over a 12\,000\,km$^2$ block;
  \item the main band's split between its components;
  \item about two main-band cycles that the passes moved over a fast trunk.
\end{itemize} What its screen still contains near a rupture is the
filter's (next subsection), and that is the same with and without the rupture.

The products add a fourth, which no split-spectrum form avoids: the main band's own integers. On D141 the PR GUNW's
main band is 25 cycles off on the cut-off block and 3--9 cycles off north of the trace. Those integers alone make our
estimate leak $-0.061$. On the Venezuela tracks the products' main band agrees with ours, and the recipe on our cells
accounts for 25--52\% of the products' leak beyond ours. Matching the products' side-band footprint changes that by
at most 0.004. The rest is in steps our matched cells do not reproduce.

\subsection{What the reordering does not remove}
\label{sec:tie}
\paragraph{The main band's integers.} The reordered estimate needs $\phiA$ unwrapped correctly, and so does the
deformation it serves. A cycle of $\phiA$ moves the raw estimate by $2\pi c_0 = 3.21$\,rad, about 62\,mm of LOS,
exactly as in the reformulated form. On the D141 pair, where SNAPHU's default and ``defo'' costs chose different
integers in a few patches, the raw estimate moved by exactly 3.21\,rad per cycle. A block that the cut isolates needs
its own integer. The range offsets give one: $m = (4\pi/\lambda)\,\rho_\mathrm{off} - \phiA$ holds no deformation, only
the block's cycles plus twice the ionosphere, because offsets see the ionosphere as a delay and phase as an advance.
Combining $m$ with the screen does not help: $m + 2\hat\phi_{io}$ moves only 0.022 cycles per cycle of the tie
\citep[Sec.~III-D makes the same point about phase--group-delay differences]{gomba2016toward}, and the bands' absolute
(``delta-$k$'') phase \citep{bovenga2013mca} is as much ionosphere as integer. The screen's own continuity across the
gap does settle it, to about a tenth of a cycle, and finds the pipelines' offset ties a cycle short on three of the
four Panama pairs (Section~\ref{sec:sbtie}). The offsets were not wrong; the pipelines' reading of them was. They
compared separate planes, fitted 6--12\,km from the trace on each side, at the trace. Read with the side band's
estimator, the same offsets agree with it to a third of a cycle. Across an Antarctic shear margin the non-dispersive phase steps by 74--188\,rad. There
the side band and the range offsets agree to 0.05 cycles, and both find the operational velocity-map fix one or two
cycles short (Section~\ref{sec:antarctica}). Far from the trace a wrong tie shifts the block's corrected LOS by a constant
59\,mm, which a per-block offset in the inversion absorbs. Near the trace the filter fills the excluded zone from both
sides and mixes them, leaving a fake step of 5--6\,cm per cycle that decays over $\sim$8\,km
(Fig.~\ref{fig:tie}d). Filtering each block separately would remove the mixing at the cost of the north side's sparse
near-trace cells.

\paragraph{The cells, the fill and the filter.} Near a rupture three parts of the estimator add error that has nothing
to do with the slip (Section~\ref{sec:nearfield}). The cells sum steep fringes differently in the two bands; the fill
of the excluded zone is a plane wider than the zone; and ISCE3's weighted mean lags the ionosphere's gradient where
the damage zone lowers the weights. On the synthetics each can be removed, by taking the main band out of every pixel,
by a fill matched to the gap and by a weighted local plane as the final step, and together they bring the error within
2\,km of the trace from 16 to 2\,mm. On real data they matter less. Real ionosphere is smoother at 2--10\,km than our
synthetics, and the fills differ by about a millimetre in a 2\,km fault zone. On D141 the plane moves the screen by at
most 3\,mm.

\paragraph{The main band's local integers.} What does matter on D141 is the main band near the trace. Its cut and uncut
unwrappings disagree on 18--78\% of the cells within 6\,km north of the trace, at 3.21\,rad of raw estimate per cycle.
Every recipe that follows the near-trace cells more closely takes those disagreements into the screen: up to 16\,mm
for the narrow fills, against 2.4\,mm or less for the production screen, whose 10\,km fill averages them away. The
production recipe's wide fill and its 2\,km exclusion are therefore the robust choice until the integers near a trace
can be checked. The side band could check them patch by patch, as it ties a cut-off block.

\paragraph{Strong ionosphere.} $r = k\,\phi_{io}$ wraps once $\phi_{io}$ changes by $2\pi/|k| = 73$\,rad. Unwrapping it
fails only if it changes by $\pi$ between neighbouring cells, about 36\,rad of ionosphere per 240\,m cell, or six cycles
of $\phiA$ per cell. The main band fails long before that. On the steepest of our screens (D126), $r$ spans about
10\,rad over the scene.

\paragraph{What it does not address.} Azimuth streaks from ionospheric scintillation, the registration errors that
ionospheric azimuth shifts cause \citep{wegmuller2024palsar3}, and the higher-order terms some sporadic-E cases need
\citep{furuya2017midlatitude} are beyond this paper. The reordering changes only what is unwrapped. It is compatible
with any filter, with the Bayesian combination of \citet{gomba2017bayesian}, and with time-series estimation
\citep{fattahi2017insar,piromthong2026ionospheric}.

\subsection{A one-line change for the operational products}
ISCE3 forms the double difference as the main-band interferogram times the conjugate of the side band's, and unwraps
it. Multiplying that product by $\exp\!\left(i\,\frac{\Delta f}{\fA}\,\phiA\right)$, using the main band's unwrapped
phase on the side band's grid, before unwrapping, and dividing it out afterwards, turns \eqref{eq:reform} into
\eqref{eq:ours} with no other change. With that, the correction passes have nothing left to correct near a rupture
and can be switched off there. The main band's own unwrapping across the rupture still matters, at 3.21\,rad of
raw estimate per cycle. The PR GUNW's unwrapped phase is 32 cycles off across the D141 rupture, because SNAPHU went
around its end, so a product that adopted the reordering would still need the rupture cut, or another way of
fixing that integer.
